\section{Combinatorics}
\subsection{Burnside's Lemma / Polya Enumeration}
\begin{minted}{c++}
using Permutation = vector<int>;
void operator*=(Permutation& p, Permutation
const& q) {
  Permutation copy = p;
  for (int i = 0; i < p.size(); i++)
    p[i] = copy[q[i]];
}
int count_cycles(Permutation p) {
  int cnt = 0;
  for (int i = 0; i < p.size(); i++) {
    if (p[i] != -1) {
      cnt++;
      for (int j = i; p[j] != -1;) {
        int next = p[j];
        p[j] = -1;
        j = next;
      }
    }
  }
  return cnt;
}
int solve(int n, int m) {
  Permutation p(n*m), p1(n*m), p2(n*m),
  p3(n*m);
  for (int i = 0; i < n*m; i++) {
    p[i] = i;
    p1[i] = (i % n + 1) % n + i / n * n;
    p2[i] = (i / n + 1) % m * n + i % n;
    p3[i] = (m - 1 - i / n) * n + (n - 1 - i % n);
  }
  set<Permutation> s;
  for (int i1 = 0; i1 < n; i1++) {
    for (int i2 = 0; i2 < m; i2++) {
      for (int i3 = 0; i3 < 2; i3++) {
        s.insert(p);
        p *= p3;
      }
      p *= p2;
    }
    p *= p1;
  }
  int sum = 0;
  for (Permutation const& p : s) {
    sum += 1 << count_cycles(p);
  }
  return sum / s.size();
}
\end{minted}
\subsection{Exclusion DP}
There is only difference to calculate $f[i]$ for different problems,

for ordered pair means $(a, b)$ and $(b, a)$ are not same then $f[i] = d[i] \cdot d[i]$;

for unordered pair means $(a, b)$ and $(b, a)$ are same then

-for repeative $\to f[i] = (d[i]+1)C2 \to d[i] \cdot (d[i]+1)/2$;

-for non-repeative $\to f[i] = d[i]C2 \to d[i] \cdot (d[i]-1)/2$;

for $k$-tuples:

-for ordered tuples : $f[i] = (d[i])^k$;

-for unordered tuples :

--for non-repeative $\to f[i] = d[i]Ck$

--for repeative $\to f[i] = (d[i]+k-1)Ck$

here $d[i]$ = number of elements divisible by $i$.

and $C$ denotes combination
\subsection{Bishop Placement}
\begin{minted}{c++}
int bishop_placements(int N, int K)
{
  if (K > 2 * N - 1)
    return 0;
  vector<vector<int>> D(N * 2, vector<int>(K +
  1));
  for (int i = 0; i < N * 2; ++i)
    D[i][0] = 1;
  D[1][1] = 1;
  for (int i = 2; i < N * 2; ++i)
    for (int j = 1; j <= K; ++j)
      D[i][j] = D[i-2][j] + D[i-2][j-1] *
      (squares(i) - j + 1);
  int ans = 0;
  for (int i = 0; i <= K; ++i)
    ans += D[N*2-1][i] * D[N*2-2][K-i];
  return ans;
}
\end{minted}
\subsection{Handling Precision (EPS)}
\begin{minted}{c++}
const double EPS = 1e-9;
bool equal(double a, double b) { return abs(a - b)
< EPS; }
bool lessThan(double a, double b) { return a < b - EPS; }
bool greaterThan(double a, double b) { return a >
b + EPS; }
\end{minted}
